\section{Recomendaciones}
\label{sec:recomendaciones}

\begin{table}[htbp]
    \centering
    \caption{Recomendaciones derivadas del an\'{a}lisis de la l\'{i}nea SIM-015.}
    \label{tab:recomendaciones}
    \setcounter{itemcount}{0}
    \begin{tabularx}{\textwidth}{@{}NlX@{}}
        \toprule
        \multicolumn{1}{c}{\textbf{Item}} & \textbf{Categoria} & \textbf{Recomendacion} \\
        \midrule
        & Construcci\'{o}n & Construir la soporter\'{i}a estrictamente conforme al isom\'{e}trico PCF114: los soportes r\'{i}gidos verticales (+Y), las gu\'{i}as y las anclas definidos en el modelo son parte integral del comportamiento verificado; cualquier cambio de posici\'{o}n o tipo de soporte invalida el presente an\'{a}lisis. \\
        & Construcci\'{o}n & Respetar la especificaci\'{o}n de espesores del isom\'{e}trico: el tramo de 3~NPS es Sch 40S y el resto de la l\'{i}nea es Sch 10S; la sustituci\'{o}n de c\'{e}dula modifica pesos y esfuerzos y requiere revalidar el modelo. \\
        & Ingenier\'{i}a & Remitir las cargas de las anclas de los nodos 140, 320 y 510 (secci\'{o}n~\ref{sec:resultados}) a los responsables de los equipos conectados para confirmar la aceptaci\'{o}n de las cargas en sus boquillas. \\
        & Documentaci\'{o}n & Conservar el job CAESAR P2603-PR-SIM-015 y el PCF corregido como registro del an\'{a}lisis; ante cualquier modificaci\'{o}n de la l\'{i}nea (di\'{a}metros, recorrido, condiciones de operaci\'{o}n), re-ejecutar el modelo y re-emitir este informe con una nueva revisi\'{o}n. \\
        \bottomrule
    \end{tabularx}
\end{table}
